% Adapted from the author's public Erdos Paper IV v1.23 series_template.tex.
% Layout provenance is recorded in 02_validation; the Markdown is semantic authority.
\documentclass[spanish,11pt,a4paper]{article}
\usepackage[margin=1in]{geometry}
\usepackage{fontspec}
\setmainfont{lmroman10-regular.otf}[BoldFont=lmroman10-bold.otf,ItalicFont=lmroman10-italic.otf,BoldItalicFont=lmroman10-bolditalic.otf]
\setsansfont{lmsans10-regular.otf}[BoldFont=lmsans10-bold.otf]
\setmonofont{lmmono10-regular.otf}[Scale=MatchLowercase]
\usepackage{amsmath,amssymb}
\providecommand{\mathscr}{\mathcal}
\usepackage[spanish,es-noshorthands,es-tabla]{babel}
\usepackage{microtype,parskip}
\usepackage{longtable,booktabs,array,calc}
\usepackage{xcolor,xurl}
\usepackage{hyperref,bookmark}
\usepackage{etoolbox,needspace}
\setcounter{secnumdepth}{-\maxdimen}
\setlength{\emergencystretch}{3em}
\providecommand{\tightlist}{\setlength{\itemsep}{0pt}\setlength{\parskip}{0pt}}
\newcounter{none}
\AtBeginEnvironment{longtable}{\Needspace{8\baselineskip}\small}
\pretocmd{\section}{\Needspace{5\baselineskip}}{}{}
\pretocmd{\subsection}{\Needspace{7\baselineskip}}{}{}
\pretocmd{\subsubsection}{\Needspace{6\baselineskip}}{}{}
\urlstyle{same}

\usepackage{graphicx}
\providecommand{\pandocbounded}[1]{#1}

\usepackage{float}

\clubpenalty=10000\widowpenalty=10000\displaywidowpenalty=10000
\usepackage{tikz}
\AtBeginEnvironment{longtable}{\Needspace{16\baselineskip}}
\hypersetup{pdftitle={Companion formal de Conteo de diseños de Steiner con perfiles de color prescritos},pdfauthor={Juan Pablo Traverso Gianini},pdflang={es},colorlinks=true,linkcolor=blue,urlcolor=blue,citecolor=blue}
\title{Companion formal de Conteo de diseños de Steiner con perfiles de color prescritos}
\author{Juan Pablo Traverso Gianini\\\small Investigador independiente, Santiago, Chile\\\small \href{mailto:jtraverso@gmail.com}{jtraverso@gmail.com}\\\small \href{https://orcid.org/0009-0003-6068-4096}{ORCID: 0009-0003-6068-4096}}\date{}\begin{document}\raggedbottom\maketitle
\begin{center}
\textbf{Paper III. Entropía combinatoria y conteo.}\\
\textbf{Preprint 1.0 --- 24 de noviembre de 2026.}\\
DOI: \url{https://doi.org/10.5281/zenodo.24501236}.\\
\href{https://github.com/jtraverso/latinon-entropy-ldp/tree/main/preprints/STEINER-PROFILES/05_formalization/lean}{Fuentes Lean en el repositorio de la serie}\end{center}
\subsection{1. Enunciados principales}\label{enunciados-principales}

El Teorema 1.1 corresponde a \path{KqR75.PhysicalCompletion.eventually_designFibre_bilateral} en \path{R75.lean}. La coloración está fijada y los puntos están etiquetados. La cuota es entera y satisface las ecuaciones de pares sobre la coloración real. La constante de error se fija antes del orden, la coloración y el histograma. Para las cuotas interiores se obtienen la existencia y ambas cotas logarítmicas. La declaración \path{KqR75.UpperCounting.eventually_card_designFibre_upper_entropy} estima directamente la cardinalidad y por eso admite también fibras vacías.

El Corolario 10.1 corresponde a \path{KqR75.ProfileLDP.designProfileLDP_probability} en \path{DesignProfileLDP.lean}. Las proporciones de las clases convergen a límites positivos a lo largo de los órdenes admisibles. Para un cerrado, el teorema admite cualquier número estrictamente inferior a la tasa en todos sus puntos factibles. Para un abierto se fija un punto factible y una tolerancia positiva. Las cotas exponenciales se refieren en ambos casos al cociente entre el número de diseños en el suceso y el número total del mismo orden. La tercera conclusión es la compacidad de los subniveles factibles.

\subsection{2. Correspondencia con las demostraciones}\label{correspondencia-con-las-demostraciones}

\textbf{Secciones 1--2.} Los márgenes se definen en \path{ExactPairMargins.lean}. La inversa sobre el completo y su norma se estudian en \path{CompleteJohnson.lean}. El Lema 2.1 aplica \path{solve_defect} de \path{HostCorrection.lean} al grafo con aristas retiradas. La corrección simultánea de cargas e histograma se encuentra en \path{JointCorrection.lean}.

\textbf{Sección 3.} El Lema 3.1 se apoya en \path{SlotPairMargins.affine_spans_kernel}. La realización de la celda sobre los vértices del diseño se verifica en \path{PhysicalAffineCell.lean}. Las declaraciones \path{pair_covered_once} y \path{switch_cell} permiten cambiar el histograma conservando la cobertura.

\textbf{Secciones 4--5.} La desigualdad para un bloque está en \path{SignedBlock.lean}. Tras sumar sobre los bloques disponibles se obtiene la deriva física \path{energy_drift_rate} de \path{PhysicalDrift.lean}. Primero se controla la proporción de peso de un tipo y después su número de selecciones. El resultado hasta el horizonte es \path{eventually_count} en \path{PhysicalHorizonCounts.lean}.

\textbf{Secciones 6--7.} La cobertura de las aristas restantes utiliza \path{eventually_cover_physical} en \path{PhysicalCover.lean}. Para el Lema 6.1, \path{OmniBridge.exists_omni} aplica el resultado importado \path{RefinedAbsorption.Original2024.theorem_1_7_two_le}; \path{eventually_absorption_budget} y \path{physical_completion} dan el grado y la cobertura que necesita el artículo. El Lema 7.1 reúne los prefijos buenos con \path{corrected_completion} de \path{CorrectedCompletion.lean}.

\textbf{Secciones 8--9.} La recuperación de los bloques elegidos y la cota factorial para la multiplicidad de sus órdenes son \path{recover_picked} y \path{order_fibre_card} en \path{PrefixRecovery.lean}. La cota superior pasa por \path{UpperWeightedReveal.lean}, \path{UpperExceptionalWeight.lean} y \path{UpperEntropyRate.lean}, y admite pesos de tipo nulos.

\textbf{Sección 10.} \path{KqR75.SlotPairMargins.round_quotas} redondea las cuotas conservando sus márgenes. \path{ProfileTypeLaw.lean} compara las leyes finita y límite. Los controles de conteo, tasa y topología están en \path{DesignProfileCounts.lean}, \path{DesignProfileRates.lean} y \path{ProfileLDPPrinciple.lean}. La equivalencia \path{orderedMargins_eq_iff} identifica la condición de pares del manuscrito con la de la declaración probabilística final. El ejemplo gráfico y la deducción elemental mediante Pinsker ilustran ese corolario en el texto.

\subsection{3. Fuentes y verificación}\label{fuentes-y-verificaciuxf3n}

El archivo \texttt{steiner-profiles-lean.zip} contiene los 162 módulos locales, la configuración de Lake, las revisiones fijadas, los verificadores y sus resultados. La entrada canónica es \path{Reconstructed} y la versión de Lean es 4.35.0. El kernel cargó 59 importaciones externas desde las dependencias compiladas y verificó 1.997 declaraciones del proyecto. El catálogo imprimió 1.750 registros de teoremas y sus 1.750 pruebas. Hay 1.741 nombres distintos, incluidos lemas privados y teoremas generados. Para cada registro se contrastaron los axiomas con \path{propext}, \path{Classical.choice} y \path{Quot.sound}.

Las firmas y los axiomas están en \texttt{verification/THEOREMS.txt}. Los términos de prueba, identificados por módulo y nombre, están en \texttt{verification/PROOFS.txt} dentro del ZIP. El archivo comprimido distribuye el catálogo completo de pruebas. Su suma de verificación figura en el \path{SHA256SUMS} adjunto. Para repetir la comprobación, descomprima el archivo y ejecute desde el directorio \path{lean}, con acceso a las dependencias fijadas y un directorio de salida nuevo:

\begin{verbatim}
lake build
lake env python ../verification/recheck.py ../recheck-output
\end{verbatim}

\end{document}
